\subsection{Loop-disentanglable spaces}

DEFINITION 2. --- A topological space B is said to be loop-disentanglable if it is locally path-connected and if every point b of B has a neighbourhood V such that the homomorphism from \(\pi_1(V,b)\) into \(\pi_1(B,b)\) deduced from the canonical injection has as its image the identity element.

A locally path-connected topological space is loop-disentanglable if and only if each of its connected components is.

Let B be a locally path-connected topological space. Suppose that every point of B has a neighbourhood V which is loop-disentanglable. Then B is loop-disentanglable.

Every locally path-connected and simply path-connected topological space is loop-disentanglable. This is in particular the case for a locally path-connected space homotopic to a point.

Remarks. --- 1) There exist topological spaces B, connected and locally path-connected, such that certain points a of B, but not all, have a neighbourhood V such that the homomorphism from \(\pi_{1}(V,a)\) into \(\pi_{1}(B,a)\) is trivial (III, p. 336, exerc. 6).

\begin{enumerate}
\def\labelenumi{\arabic{enumi})}
\setcounter{enumi}{1}
\tightlist
\item
  Let B be a loop-disentanglable space; every operation of the groupoid \(\varpi(\mathrm{B})\) on a B-space is without local monodromy (cf. III, p. 313, remark). In particular, for a map \(p: E \rightarrow B\) to make E a covering of B, it is necessary and sufficient that it be étale, separated, and that it satisfy the path-lifting property (III, p. 315, corollary 3).
\end{enumerate}

PROPOSITION 1. --- The product space of a finite family of loop-disentanglable spaces is loop-disentanglable.

It suffices to prove that, if A and B are two loop-disentanglable spaces, then their product A × B is also. Under these conditions, the space A × B is indeed locally path-connected (III, p. 261, prop. 9). Let (a, b) be a point of A × B. By hypothesis there exists a neighbourhood U of a (resp. a neighbourhood V of b) such that the image of the homomorphism \(i_* \colon \pi_1(\mathrm{U}, a) \to \pi_1(\mathrm{A}, a)\) deduced from the canonical injection \(i \colon \mathrm{U} \to \mathrm{A}\) (resp. of the homomorphism \(j_* \colon \pi_1(\mathrm{V}, b) \to \pi_1(\mathrm{B}, b)\) deduced from the canonical injection \(j \colon \mathrm{V} \to \mathrm{B}\) ) is reduced to the identity element. Then U × V is a neighbourhood of (a, b) in A × B. The homomorphism \(\pi_1(\mathrm{U} \times \mathrm{V}, (a, b)) \to \pi_1(\mathrm{A} \times \mathrm{B}, (a, b))\) is identified with the homomorphism \((i_*, j_*)\) (III, p. 297, corollary of prop. 4). Its image is therefore reduced to the identity element. This proves that A × B is loop-disentanglable.

PROPOSITION 2. --- a) Every covering of a loop-disentanglable space is loop-disentanglable.

\begin{enumerate}
\def\labelenumi{\alph{enumi})}
\setcounter{enumi}{1}
\item
  Conversely, let \(p \colon \mathrm{E} \to \mathrm{B}\) be an étale and surjective map and suppose that \(\mathrm{E}\) is a loop-disentanglable space. Then \(\mathrm{B}\) is loop-disentanglable.
\item
  Let B be a loop-disentanglable topological space and let \((\mathrm{E}, p)\) be a covering of B. The space E is then locally path-connected. Let x be a point of E, write \(b = p(x)\) , and let V be a neighbourhood of b in B such that the canonical homomorphism \(\pi_{1}(\mathrm{V}, b) \to \pi_{1}(\mathrm{B}, b)\) is trivial. Write \(i: V \to B\) and \(j: p^{-1}(V) \to E\) for the canonical injections. By hypothesis, the image of the homomorphism \(\pi_{1}(i, b)\) is reduced to the identity element. Since \(p \circ i = j \circ p\) and \(\pi_{1}(p, x)\) is injective, the image of the homomorphism \(\pi_{1}(j, x)\) is reduced to the identity element. This proves that E is loop-disentanglable.
\item
  The space B is locally path-connected. Let b be a point of B and let x be a point of \(E_{b}\) . Since p is étale, there exists a neighbourhood U of x in E such that p induces a homeomorphism of U onto \(p(\mathrm{U})\) . Let V be a neighbourhood of x in E contained in U such that the homomorphism \(\pi_{1}(j,x)\) is trivial, where \(j:V\to E\) denotes the canonical injection. Write \(q:V\to p(V)\) for the map deduced from p by passing to subspaces; it is a homeomorphism. Also write i for the canonical injection of \(p(\mathrm{V})\) into B. We have \(p\circ j=i\circ q\) , therefore \(\pi_{1}(i,b)\circ\pi_{1}(q,x)=\pi_{1}(p,x)\circ\pi_{1}(j,x)\) is the trivial homomorphism. Since the homomorphism \(\pi_{1}(q,x)\) is an isomorphism, the homomorphism \(\pi_{1}(i,b)\) is trivial. It follows that the space B is loop-disentanglable.
\end{enumerate}

PROPOSITION 3. --- Let B be a loop-disentanglable topological space. Let (Y, q) be a covering of B and let (Z, p) be a covering of Y. The topological space Z, equipped with the map \(q \circ p\) , is then a covering of B.

Indeed, the map \(q \circ p\) is étale (I, p. 29, prop. 6) and separated (I, p. 25, prop. 2). By remark 2 above, it therefore suffices to show that it satisfies the path-lifting property. Let \(z\) be a point of Z and let \(c \colon \mathbf{I} \to \mathbf{B}\) be a path in B with origin \(q(p(z))\) . There exists a path \(c'\) with source \(p(z)\) in Y which lifts \(c\) , since Y is a covering of B (III, p. 302, cor. 2). Since Z is a covering of Y, there exists a path \(c''\) with source \(z\) in Z which lifts \(c'\) ; the path \(c''\) lifts \(c\) . This proves the proposition.

